\scanpage{549}

et on suppose que $\unsure{\mathcal{G}|\mathfrak{G}}^+\ne\unsure{\mathcal{G}|\mathfrak{G}}$, donc
\[
\tag{3}
\begin{gathered}
\unsure{\mathcal{G}|\mathfrak{G}}\simeq\{1,\tau'\}\mathbin{.}\unsure{\mathcal{G}|\mathfrak{G}}^+\\[-2pt]
(\tfrac12\ \text{direct})
\end{gathered}
\]
et on a une suite exacte canonique
\[
\tag{4}
\diagram{\begin{tikzcd}[column sep=large]
1\arrow[r]&\unsure{\mathcal{G}|\mathfrak{G}}^+\arrow[r]&\unsure{\mathcal{G}|\mathfrak{G}}\arrow[r,"\varepsilon_{\mathcal{G}}"]&\{\pm1\}\arrow[r]&1
\end{tikzcd}}
\]
On désigne aussi par $\nu_{\mathcal{G}}$ le composé
\[
\tag{5}
\diagram{\begin{tikzcd}[column sep=large,row sep=small]
\unsure{\mathcal{G}|\mathfrak{G}}\arrow[r,"\varepsilon_{\mathcal{G}}"]\arrow[rr,bend left=25,"\nu_{\mathcal{G}}"]&\{\pm1\}&\mathbb{Z}/2\mathbb{Z}\arrow[l,"\sim"']\\
&(-1)^\nu&\nu\arrow[l,mapsto]
\end{tikzcd}}\ .
\]

De façon générale, on \unsure{désignera} (par abus de langage) par une $\overline{\mathrm m}$ lettre des éléments de $\mathrm{Gl}(2,\mathbb{Z})$ --- tels $\varepsilon_0,\varepsilon_1,\tau=\tau_\infty,\tau'=\tau'_\infty$, $\sigma=\sigma_\infty,\rho=\rho_\infty$ etc --- et leurs images dans $\unsure{\mathcal{G}|\mathfrak{G}}$.

\margin{Les groupes $Z_{\rho,\sigma}$, $\gamma_{\rho,\sigma}$}

$2^\circ)$ Soit
\[
\tag{6}
\begin{gathered}
\unsure{Z_{\rho,\sigma}|\mathbb{Z}_{\rho,\sigma}}\ \unsure{\subset|=}\ \operatorname{Aut}(\unsure{\mathcal{G}|\mathfrak{G}}^+)\times\widehat{\mathbb{Z}}^*\times\{\pm1\}\times\{\pm1\}\\[-2pt]
\Vert\ \mathrm{dfn}\\[-2pt]
\left\{(u,\mu,\varepsilon,\varepsilon')\ \middle|\
\begin{aligned}
u(\varepsilon_0)&=\varepsilon_0^\mu\\
u(\rho)&\ \text{conjugué dans }\unsure{\mathcal{G}|\mathfrak{G}}^+\text{ à }\rho^{\varepsilon'}\\
u(\sigma)&\ \rule{4.1cm}{0.3pt}\ \text{à }\sigma^\varepsilon
\end{aligned}\right\}
\end{gathered}
\]
C'est un sous-groupe du groupe produit.

Posons
\[
\tag{7}
\left\{\begin{aligned}
L_\rho&=\text{sous-groupe de }\unsure{\mathcal{G}|\mathfrak{G}}\text{ engendré par }\rho\\
L_\sigma&=\rule{7cm}{0.3pt}\ \sigma
\end{aligned}\right.
\]
de sorte qu'on a des \unsure{épimorphismes}
\[
\tag{8}
\left\{\diagram{\begin{tikzcd}[column sep=large,row sep=small]
\mathbb{Z}/6\mathbb{Z}\arrow[r,"\text{épi}"]&L_\rho&n\arrow[r,mapsto]&\rho^n\\
\mathbb{Z}/4\mathbb{Z}\arrow[r]&L_\sigma&n\arrow[r,mapsto]&\unsure{\sigma^n}
\end{tikzcd}}\right.\ ,
\]
\mspage{548}

{\small
\begin{center}
548\\
\texttt{https://grothendieck.umontpellier.fr}\\
Université de Montpellier, UMR CNRS 5149, Case courrier 051, Place Eugène Bataillon\\
34095 Montpellier cedex 5 - France
\end{center}
}

\scanpage{550}
\mspage{578}

et \deleted{des} \inserted{un diagramme \unsure{commut$^{\mathrm{if}}$|commut'}} \unsure{d'}homomorphismes
\[
\tag{9}
\diagram{\begin{tikzcd}[column sep=large,row sep=small]
&L_\rho\arrow[dr,hook]&\\
\{\pm1\}\arrow[ur]\arrow[dr]&&\unsure{\mathcal{G}|\mathfrak{G}}^+\\
&L_\sigma\arrow[ur,hook]&
\end{tikzcd}}
\]

--- Si $\unsure{\mathcal{G}|\mathfrak{G}}$ est ``assez gros'', on aura d'ailleurs que $\{\pm1\}\hookrightarrow\unsure{\mathcal{G}|\mathfrak{G}}^+$ i.e. ``$-1$'' $\ne1$ dans $\unsure{\mathcal{G}|\mathfrak{G}}^+$, et même
\[
L_\rho\cap L_\sigma=\{\pm1\}\ .
\]

Les conditions (6) impliquent donc que
\[
\tag{10}
\left\{\begin{aligned}
u(L_\rho)&\ \inserted{\unsure{\mathcal{G}|\mathfrak{G}}^+\text{-}}\text{conjugué à }L_\rho\\
u(L_\sigma)&\ \unsure{\mathcal{G}|\mathfrak{G}}^+\text{-conjugué à }L_\sigma
\end{aligned}\right.
\]

Si $\unsure{\mathcal{G}|\mathfrak{G}}^+$ est ``assez gros'' pour que $\rho$ et $\rho^{-1}$, $\sigma$ et $\sigma^{-1}$ n'y soient pas conjugués, alors un élément $(u,\mu,\varepsilon,\varepsilon')$ de $Z_{\rho,\sigma}$ est connu par la seule connaissance de $(u,\mu)$. Si d'autre part l'hom
\[
\tag{11}
\diagram{\begin{tikzcd}[column sep=large,row sep=small]
\widehat{\mathbb{Z}}\arrow[r]&\unsure{\mathcal{G}|\mathfrak{G}}^+\\
n\arrow[r,mapsto]&\varepsilon_0^n
\end{tikzcd}}
\]
est injectif, alors $\mu$ est connu quand on connaît $u$. Donc si $\unsure{\mathcal{G}|\mathfrak{G}}^+$ est assez gros pour satisfaire aux trois conditions envisagées (p.ex. pour qu'on ait un ``hom modulaire''
\[
\tag{12}
\diagram{\begin{tikzcd}[column sep=large]
\unsure{\mathcal{G}|\mathfrak{G}}\arrow[r]&\mathrm{Gl}(2,\widehat{\mathbb{Z}})
\end{tikzcd}}
\]
(transformant $\rho,\sigma,\tau'$ en les éléments de $\overline{\mathrm m}$ noms), alors un élément de $Z_{\rho,\sigma}$ est connu par la seule
\mspage{549}

{\small
\begin{center}
549\\
\texttt{https://grothendieck.umontpellier.fr}\\
Université de Montpellier, UMR CNRS 5149, Case courrier 051, Place Eugène Bataillon\\
34095 Montpellier cedex 5 - France
\end{center}
}

\scanpage{551}

connaissance de $u\in\operatorname{Aut}(\unsure{\mathcal{G}|\mathfrak{G}}^+)$.

On a un hom
\[
\tag{13}
\diagram{\begin{tikzcd}[column sep=large,row sep=small]
L_0\arrow[r,"i_{L_0,Z}"]&\unsure{Z_{\rho,\sigma}}\\
\ell\arrow[r,mapsto]&(\operatorname{int}(\ell),1,1,1)
\end{tikzcd}}
\]
dont l'image est \unsure{invariante}, le noyau n'est autre que $L_0\cap\operatorname{Cent}(\unsure{\mathcal{G}|\mathfrak{G}}^+)$. Notons que l'on a
\[
\tag{14}
L_0\cap\operatorname{Cent}(\unsure{\mathcal{G}|\mathfrak{G}}^+)\subset L_0\cap\operatorname{Cent}_{\mathcal{G}^+}(\rho)\subset L_0\cap L_1
\]
car si $\ell\in L_0$ et $\rho$ commutent, \inserted{i.e. \uline{si $\ell$}$=\rho(\ell)$,} comme \deleted{$\ell=\rho(\ell)$} $\rho(L_0)=L_1$, on aura $\ell\in L_1$. Donc si $L_0\cap L_1=\{1\}$ (ce qui est le cas si (12) existe), \deleted{alors} et $\overline{\mathrm m}$ s'il existe seulement
\[
\tag{12'}
\diagram{\begin{tikzcd}[column sep=large]
\unsure{\mathcal{G}|\mathfrak{G}}\arrow[r]&\mathrm{Gl}(2,\widehat{\mathbb{Z}})'
\end{tikzcd}}\qquad\text{modulaire},
\]
alors on aura
\[
\tag{15}
L_0\cap\operatorname{Cent}(\unsure{\mathcal{G}|\mathfrak{G}}^+)=\{1\},
\]
et (13) sera injectif. Nous allons le supposer, pour fixer les idées.

On pose
\[
\tag{16}
\gamma_{\rho,\sigma}=Z_{\rho,\sigma}/\operatorname{Im}L_0,
\]
de sorte qu'on aura une suite exacte canonique
\[
\tag{17}
\diagram{\begin{tikzcd}[column sep=large]
1\arrow[r]&L_0\arrow[r,"i_{L_0,Z}"]&Z_{\rho,\sigma}\arrow[r,"\delta_Z"]&\gamma_{\rho,\sigma}\arrow[r]&1
\end{tikzcd}}\ .
\]

Les hom. canoniques de projection $(u,\mu,\varepsilon,\varepsilon')\mapsto\mu$, $\unsure{\varepsilon},\varepsilon'$ se factorisent par $\gamma_{\rho,\sigma}$ et fournissent des homomorphismes canoniques
\mspage{550}

{\small
\begin{center}
550\\
\texttt{https://grothendieck.umontpellier.fr}\\
Université de Montpellier, UMR CNRS 5149, Case courrier 051, Place Eugène Bataillon\\
34095 Montpellier cedex 5 - France
\end{center}
}

\scanpage{552}
\mspage{579}

\[
\tag{18}
\left\{\diagram{\begin{tikzcd}[column sep=large,row sep=small]
\gamma_{\rho,\sigma}\arrow[r,"\chi_\gamma"]&\widehat{\mathbb{Z}}^*\\
\gamma_{\rho,\sigma}\arrow[r,"\varepsilon_\gamma"]&\{\pm1\}\\
\gamma_{\rho,\sigma}\arrow[r,"\varepsilon'_\gamma"]&\{\pm1\}
\end{tikzcd}}\right.
\quad\text{soit}\quad
\diagram{\begin{tikzcd}[column sep=huge]
\gamma_{\rho,\sigma}\arrow[r,"(\chi_\gamma,\varepsilon_\gamma,\varepsilon'_\gamma)"]&\widehat{\mathbb{Z}}^*\times\{\pm1\}\times\{\pm1\}
\end{tikzcd}}
\]
les composés de ceux-ci avec $Z_{\rho,\sigma}\to\gamma_{\rho,\sigma}$ sont notés
\[
\tag{19}
\chi_Z=\chi_\gamma\circ\delta_Z,\quad
\unsure{\varepsilon_Z|\xi_Z}=\varepsilon_\gamma\circ\delta_Z,\quad
\varepsilon'_Z=\varepsilon'_\gamma\circ\delta_Z,
\]
de sorte que
\[
\tag{20}
\chi_Z(u,\mu,\varepsilon,\varepsilon')=\mu,\quad
\unsure{\varepsilon_Z|\xi_Z}(u,\mu,\varepsilon,\varepsilon')=\varepsilon,\quad
\varepsilon'_Z(u,\mu,\varepsilon,\varepsilon')=\varepsilon'
\]

On aurait envie de désigner par $\gamma_{\rho,\sigma}^{\unsure{0|\circ}}$ le noyau de $(\varepsilon_\gamma,\varepsilon'_\gamma)$, par $\gamma_{\rho,\sigma}^{+}$ celui de $\chi_\gamma$, par $\gamma_{\rho,\sigma}^{\unsure{0|\circ}+}$ leur intersection, et d'introduire par les notations $Z_{\rho,\sigma}^{\unsure{0|\circ}}$, $Z_{\rho,\sigma}^{+}$, $Z_{\rho,\sigma}^{\unsure{0|\circ}+}$ leurs images inverses dans $Z_{\rho,\sigma}$, i.e. les noyaux \unsure{resp.$^{\mathrm t}$|resp.} de $(\varepsilon_Z,\varepsilon'_Z)$, de $\chi_Z$ et de $(\chi_Z,\varepsilon_Z,\varepsilon'_Z)$.

\margin{Les groupes $S_0G_{\rho,\sigma}$, $G_{\rho,\sigma}$}

$3^\circ)$ \deleted{Le groupe $Z_{\rho,\sigma}$ \unsure{opère},} \inserted{On a un hom. tautologique}
\[
\tag{21}
\diagram{\begin{tikzcd}[column sep=large,row sep=small]
Z_{\rho,\sigma}\arrow[r]&\operatorname{Aut}(\unsure{\mathcal{G}|\mathfrak{G}}^+)\\
(u,\mu,\varepsilon,\varepsilon')\arrow[r,mapsto]&u
\end{tikzcd}}\ ,
\]
donc $Z_{\rho,\sigma}$ \unsure{agit|opère} sur $\unsure{\mathcal{G}|\mathfrak{G}}^+$, et on peut donc considérer le produit $\tfrac12$ direct
\[
\tag{22}
\begin{gathered}
S_0G_{\rho,\sigma}=Z_{\rho,\sigma}\mathbin{.}\unsure{\mathcal{G}|\mathfrak{G}}^+\ .\\[-2pt]
(\tfrac12\text{-direct})
\end{gathered}
\]
\mspage{551}

{\small
\begin{center}
551\\
\texttt{https://grothendieck.umontpellier.fr}\\
Université de Montpellier, UMR CNRS 5149, Case courrier 051, Place Eugène Bataillon\\
34095 Montpellier cedex 5 - France
\end{center}
}

\scanpage{553}

On a un hom
\[
\tag{23}
\diagram{\begin{tikzcd}[column sep=large,row sep=small]
L_0\arrow[r,"i_{L_0,SG}"]&S_0G_{\rho,\sigma}=Z_{\rho,\sigma}\mathbin{.}\unsure{\mathcal{G}|\mathfrak{G}}^+\\
\ell\arrow[r,mapsto]&i_{L_0,Z}(\ell)\mathbin{.}\ell^{-1}
\end{tikzcd}}
\]
dont l'image est encore invariante, \deleted{\unsure{Posons|On a}} \inserted{On} \deleted{\unsure{est stable}} \inserted{\unsure{pose}}
\[
\tag{24}
G_{\rho,\sigma}=\underbrace{S_0G_{\rho,\sigma}}_{Z_{\rho,\sigma}\mathbin{.}\unsure{\mathcal{G}|\mathfrak{G}}^+}/\operatorname{Im}L_0\ \deleted{\unsure{Z_\rho|\mathcal{G}}},
\]
de sorte qu'on aura une suite exacte
\[
\tag{25}
\diagram{\begin{tikzcd}[column sep=large]
1\arrow[r]&L_0\arrow[r,no head,"i_{L_0,SG}"]&S_0G_{\rho,\sigma}\arrow[r]&G_{\rho,\sigma}\arrow[r,no head]&1
\end{tikzcd}}\ ,
\]
et \unsure{un} diagramme \inserted{\unsure{comm.}} d'homomorphismes
\[
\tag{26}
\diagram{\begin{tikzcd}[column sep=large,row sep=small]
&Z_{\rho,\sigma}\arrow[dr,hook,"i_Z"]&\\
L_0\arrow[ur,hook,"i_{L_0,Z}"]\arrow[dr,hook,"i_{L_0,\mathcal{G}}"']&&G_{\rho,\sigma}\\
&\unsure{\mathcal{G}|\mathfrak{G}}^+\arrow[ur,hook,"i_{\mathcal{G}^+}"']&
\end{tikzcd}}
\]
(a priori $\operatorname{Ker}(\unsure{\mathcal{G}|\mathfrak{G}}^+\to G_{\rho,\sigma})\simeq\operatorname{Ker}(L_0\to Z_{\rho,\sigma})\unsure{\simeq|=}L_0\cap\operatorname{Cent}(\unsure{\mathcal{G}|\mathfrak{G}}^+)$, mais par hyp. (15) il est réduit à $\{1\}$), s'insérant dans un diagramme comm. de suites exactes
\[
\tag{27}
\diagram{\begin{tikzcd}[column sep=large,row sep=large]
1\arrow[r]&L_0\arrow[r,"i_{L_0,Z}"]\arrow[d,hook,"i_{L_0,\mathcal{G}}"']&Z_{\rho,\sigma}\arrow[r,"\delta_Z"]\arrow[d,hook]&\gamma_{\rho,\sigma}\arrow[r]\arrow[d,equal]&1\\
1\arrow[r]&\unsure{\mathcal{G}|\mathfrak{G}}^+\arrow[r,"i_{\mathcal{G}^+}"]&G_{\rho,\sigma}\arrow[r,"\delta_G"]&\gamma_{\rho,\sigma}\arrow[r,no head]&1
\end{tikzcd}}\ .
\]
On posera
\[
\tag{28}
\chi_G=\chi_\gamma\circ\delta_G,\quad
\varepsilon_G=\varepsilon_\gamma\circ\delta_G,\quad
\varepsilon'_\gamma\circ\delta_G
\]
\mspage{552}

{\small
\begin{center}
552\\
\texttt{https://grothendieck.umontpellier.fr}\\
Université de Montpellier, UMR CNRS 5149, Case courrier 051, Place Eugène Bataillon\\
34095 Montpellier cedex 5 - France
\end{center}
}

\scanpage{554}
\mspage{580}

d'où un hom
\[
\tag{29}
\diagram{\begin{tikzcd}[column sep=large]
G_{\rho,\sigma} \arrow[r,"{(\chi_G,\varepsilon_G,\varepsilon'_G)}"] &
\widehat{\mathbb Z}^{*}\times\{\pm1\}\times\{\pm1\}
\end{tikzcd}},
\]
nous conduisant à introduire $G_{\rho,\sigma}^{+}$, $G_{\rho,\sigma}^{\unsure{0}}$,
$G_{\rho,\sigma}^{+0}$ comme d'habitude.

L'injection $i_{\mathcal G^+}:\unsure{\mathcal G^+|\widetilde G^+}\hookrightarrow G_{\rho,\sigma}$
peut se prolonger en un plongement \deleted{\unsure{évident}}
(noté $i_{\mathcal G}$)
\[
\tag{30}
\diagram{\begin{tikzcd}
i_{\mathcal G}:\mathcal G \arrow[r,hook] & G_{\rho,\sigma}
\end{tikzcd}},
\]
que je définis par sa valeur sur \deleted{$\unsure{\tau'}$}
\inserted{$\tau=\sigma\tau'$}, par
\[
\tag{31}
i_{\mathcal G}(\unsure{\tau|\tau'})=
\deleted{\unsure{i_{\mathcal G^+}(\tau)}}\,i_Z(\tau_Z)
\]
\margin{Si $\tau_\gamma$ désigne l'image de $\tau=\tau_G\in G_{\rho,\sigma}$
dans $\gamma_{\rho,\sigma}=G_{\rho,\sigma}/\mathcal G^+$,
$\mathcal G$ s'identifie à l'image inverse dans $G_{\rho,\sigma}$ du sous-groupe
$\{1,\tau_\gamma\}$ de $\gamma_{\rho,\sigma}$.}
\[
\tag{32}
\left\{\begin{aligned}
&\tau_Z\in Z_{\rho,\sigma}\quad\text{défini comme}\\
&\tau_Z=\left(
\underbrace{\operatorname{int}(\tau)|\mathcal G^+}_{\in\operatorname{Aut}(\mathcal G^+)},
\underset{\substack{\rotatebox{90}{$\in$}\\\widehat{\mathbb Z}^{*}}}{-1},
\underset{\substack{\rotatebox{90}{$\in$}\\\{\pm1\}}}{-1},
\underset{\substack{\rotatebox{90}{$\in$}\\\{\pm1\}}}{-1}\right).
\end{aligned}\right.
\]

Désormais nous identifions $\mathcal G$ à un sous-groupe de $G_{\rho,\sigma}$,
et nous désignons par la $\overline{\mathrm m}$ lettre des éléments de $\mathcal G$
(tels que $\tau,\tau',\rho,\sigma,\unsure{\varepsilon_0|\xi_0},\unsure{\varepsilon_1|\xi_1}$, etc)
et leurs images dans $G_{\rho,\sigma}$.
\margin{Noter que $\tau'_G\notin Z_{\rho,\sigma}$ par contre on a
\deleted{$\unsure{\tau'_G}$} $\tau_G\in Z_{\rho,\sigma}$
\deleted{\illegible{1}}.}
On aura d'ailleurs
\[
\tag{32bis}
\chi_G|\mathcal G=\varepsilon_G|\mathcal G=\varepsilon'_G|\mathcal G
=\varepsilon_{\mathcal G}:\left\{\begin{array}{l}
\text{trivial sur }\mathcal G^+\\
-1\text{ sur }\mathcal G^-=\mathcal G\setminus\mathcal G^+.
\end{array}\right.
\]

\margin{Les groupes $Z_\rho,Z_\sigma,N_\rho,N_\sigma,SN_\rho,SN_\sigma$\unsure{\dots}}
$4^\circ)$ On \deleted{définit} note que l'élément $\tau'$ de
\deleted{$\unsure{G_{\rho,\sigma}}$} $\mathcal G\subset G_{\rho,\sigma}$
normalise $L_\rho$ et $L_\sigma$, donc normalise aussi
\[
\tag{33}
\left\{\begin{aligned}
Z_\rho&\overset{\mathrm{dfn}}=\operatorname{Cent}_{G_{\rho,\sigma}}(\rho)
=\operatorname{Cent}_{G_{\rho,\sigma}}(L_\rho),\\
Z_\sigma&\overset{\mathrm{dfn}}=\operatorname{Cent}_{G_{\rho,\sigma}}(\sigma)
=\operatorname{Cent}_{G_{\rho,\sigma}}(L_\sigma)
\end{aligned}\right.
\]
\mspage{553}

{\small
\begin{center}
553\\
\texttt{https://grothendieck.umontpellier.fr}\\
Université de Montpellier, UMR CNRS 5149, Case courrier 051, Place Eugène Bataillon\\
34095 Montpellier cedex 5 - France
\end{center}
}

\scanpage{555}

donc $\{1,\tau'\}$ agit sur ces groupes, et on peut considérer
\[
\tag{34}
\left\{\begin{aligned}
N_\rho&=\{1,\tau'\}\mathbin{.}Z_\rho,\qquad\deleted{\unsure{N}}\\
N_\sigma&=\{1,\tau'\}\mathbin{.}Z_\sigma
\end{aligned}\right.
\qquad(\text{produits }\tfrac12\text{ directs}).
\]

On a des suites exactes canoniques
\[
\tag{35}
\diagram{\begin{tikzcd}
1\arrow[r]&\unsure{Z_\rho|L_\rho}\arrow[r]&N_\rho\arrow[r,"\varepsilon_\rho"]&\{\pm1\}\arrow[r]&1\\
1\arrow[r]&Z_\sigma\arrow[r]&N_\sigma\arrow[r,"\varepsilon_\sigma"]&\{\pm1\}\arrow[r]&1
\end{tikzcd}},
\]
et un diagramme commutatif d'hom. canoniques
\[
\tag{36}
\diagram{\begin{tikzcd}[row sep=small,column sep=large]
&\unsure{\mathbf D_\rho|\mathbb D_\rho}\arrow[r]&N_\rho\arrow[dr,"i_\rho"]&\\
\{1,\tau'\}\times\{\pm1\}\arrow[ur,hook]\arrow[dr,hook]&&&G_{\rho,\sigma}\\
&\mathbf D_\sigma\arrow[r]&N_\sigma\arrow[ur,"\unsure{i_\sigma}"']&
\end{tikzcd}}
\]
où
\[
\tag{37}
\left\{\begin{aligned}
\mathbf D_\rho&\simeq\mathbf D_6\subset\operatorname{Gl}(2,\mathbb Z)
\quad\text{engendré par }\rho,\tau'\\
\mathbf D_\sigma&\simeq\mathbf D_4\subset\operatorname{Gl}(2,\mathbb Z)
\quad\text{\rule{2.7cm}{0.35pt}}\ \sigma,\tau'\\
\{1,\tau'\}\times\{\pm1\}&=\mathbf D_\rho\cap\mathbf D_\sigma
\quad\text{dans }\operatorname{Gl}(2,\mathbb Z),
\end{aligned}\right.
\]
\deleted{\unsure{Les hom} $\mathbf D_\rho\longrightarrow N_\rho$}
\[
\tag{38}
\left\{\begin{aligned}
\mathbf D_\rho&=\{1,\tau'\}\mathbin{.}L_\rho
\quad(\text{du moins si }\rho\text{ d'ordre }6\text{ dans }\mathcal G)\\
\mathbf D_\sigma&=\{1,\tau'\}\mathbin{.}L_\sigma
\quad(\text{\rule{1.5cm}{0.35pt}}\ \sigma\text{ d'ordre }4\ \text{\rule{1.4cm}{0.35pt}}),
\end{aligned}\right.
\]
et les hom
\[
\diagram{\begin{tikzcd}
\mathbf D_\rho\arrow[r]&N_\rho,
\end{tikzcd}}
\qquad
\diagram{\begin{tikzcd}
\mathbf D_\sigma\arrow[r]&N_\sigma
\end{tikzcd}}
\]
sont évidents. Les hom $i_\rho,i_\sigma$ aussi, bien sûr.

Rappelons que
\[
\tag{39}
\left\{\begin{aligned}
\deleted{\unsure{i'_\rho}}i_\rho\text{ injectif}
&\Longleftrightarrow\tau'\notin Z_\rho
\Longleftrightarrow\rho^2\ne1
\Longleftrightarrow\rho\ne-1
\Longleftrightarrow\mathcal G^+\text{ est }\unsure{\deleted{\text{isom}}|\text{isom}}
\text{ à un quotient de }\mathbb Z/4\mathbb Z.\\
i_\sigma\text{ injectif}
&\Longleftrightarrow\tau'\notin Z_\sigma
\Longleftrightarrow\sigma^2\ne1\quad
\deleted{\unsure{\Longleftrightarrow\sigma\ne\pm1}}
\end{aligned}\right.
\]
\mspage{554}

{\small
\begin{center}
554\\
\texttt{https://grothendieck.umontpellier.fr}\\
Université de Montpellier, UMR CNRS 5149, Case courrier 051, Place Eugène Bataillon\\
34095 Montpellier cedex 5 - France
\end{center}
}

\scanpage{556}
\mspage{581}

Faisant opérer $G_{\rho,\sigma}$ sur $\mathcal G^+$, sous-groupe distingué,
via autom. int., on voit que tout élément de $G_{\rho,\sigma}$ transforme
$L_\rho\subset\mathcal G^+$ en un sous-groupe \inserted{$\mathcal G^+$-}conjugué,
donc est \unsure{congru} mod \deleted{$\unsure{\operatorname{Norm}_{\mathcal G^+}(L_\rho)}$}
$\mathcal G^+$ à un élément de
$\operatorname{Norm}_{\unsure{G|\mathcal G}}(L_\rho)=\unsure{\operatorname{Im}}i_\rho(N_\rho)$,
et de même pour $L_\sigma$. En d'autres termes, \inserted{posant}
\inserted{\[
\tag{40}
\begin{aligned}
N_\rho^*&\overset{\mathrm{dfn}}=\{x\in N_\rho\mid\varepsilon_\rho(x)=\varepsilon'_G(i_\rho(x))\},\qquad
N_\sigma^*\overset{\mathrm{dfn}}=\{x\in N_\sigma\mid\varepsilon_\sigma(x)=\varepsilon_G(i_\sigma(x))\}
\end{aligned}
\]}
les hom composés
\[
\tag{\unsure{40}}
\left\{\begin{array}{ll}
\diagram{\begin{tikzcd}[column sep=large]
N_\rho^*\arrow[r,"i_\rho"]&G_{\rho,\sigma}\arrow[r,"\delta_G"]&\gamma_{\rho,\sigma}
\end{tikzcd}},&\delta_\rho\overset{\mathrm{dfn}}=\delta_G\circ i_\rho\\[1.5em]
\diagram{\begin{tikzcd}[column sep=large]
N_\sigma^*\arrow[r,"i_\sigma"]&G_{\rho,\sigma}\arrow[r,"\delta_G"]&\gamma_{\rho,\sigma}
\end{tikzcd}},&\delta_\sigma\overset{\mathrm{dfn}}=\delta_G\circ i_\sigma
\end{array}\right.
\]
des épimorphismes.
\deleted{On \unsure{voit} \illegible{3}}
\deleted{\unsure{On a des propriétés}}

\margin{Si $g\in G_{\rho,\sigma}$, et
$\varepsilon=\varepsilon_G(g)$, $\varepsilon'=\varepsilon'_G(g)$, alors
\[
\begin{aligned}
\operatorname{int}(g)(\rho)&\underset{\mathcal G^+}{\sim}\rho^{\varepsilon'}
=\unsure{\tau^{\nu'}(\rho)},\\
\operatorname{int}(g)(\sigma)&\underset{\mathcal G^+}{\sim}\sigma^\varepsilon
=\unsure{\tau^\nu(\rho)|\tau^\nu(\sigma)}.
\end{aligned}
\]
\deleted{\unsure{donc}}
(où $(-1)^\nu=\varepsilon$, $(-1)^{\nu'}=\varepsilon'$), \inserted{donc}
\deleted{$\exists g\in\unsure{\mathcal G^+}$ tel que}
$\exists a=\tau^{\nu'}\mathbin{.}\unsure{a_0|ag}\in N_\rho$
(où $a_0=\unsure{\tau^\nu g}$) tel que

a) $g\equiv\unsure{a|a_0}\ (\mathcal G^+)$ et

b) $\varepsilon_\rho(a)=\unsure{\varepsilon'|\varepsilon'_G}
=\unsure{\varepsilon'_G(g)|\varepsilon_G(g)}$ ---

(et idem pour $N_\sigma$\unsure{\dots})

donc
\[
\begin{aligned}
\operatorname{Im}\bigl(N_\rho^{\deleted{\unsure{*}}(+1)}
\xrightarrow{\deleted{\unsure{\delta_\rho}}}\gamma_{\rho,\sigma}\bigr)
&\supset\gamma'_{\rho,\sigma}\deleted{\unsure{\varepsilon'}},\\
\operatorname{Im}\bigl(N_\rho^{(-1)}\longrightarrow\gamma_{\rho,\sigma}\bigr)
&\supset\gamma_{\rho,\sigma}\setminus\gamma'_{\rho,\sigma}.
\end{aligned}
\]
Mais celà ne prouve \uline{pas} que l'on ait \unsure{commutativité}!
\unsure{d'où sans doute}
}

\deleted{\unsure{Ils ne sont pas} \illegible{2} (cf (39))}
\deleted{(\unsure{41}) a) Si $\rho^2\ne1$, alors on a commutativité de}
\[
\diagram{\deleted{\begin{tikzcd}[column sep=large,row sep=small]
N_\rho\arrow[rr,"\delta_\rho"]
\arrow[dr,"\varepsilon_\rho"']&&
\gamma_{\rho,\sigma}\arrow[dl,"\varepsilon'_\gamma"]\\
\unsure{\varepsilon_\rho}&\{\pm1\}&\\
&\unsure{\{\pm1\}}&\unsure{\{\pm1\}}
\end{tikzcd}}}
\]
\deleted{d'où}
diagramme \inserted{comm.} de \deleted{$\unsure{\{\pm1\}}$} suites exactes
\[
\tag{41}
\diagram{\begin{tikzcd}[column sep=large]
1\arrow[r]&Z_\rho\arrow[r]\arrow[d,"\text{épi}"]&N_\rho^*\arrow[r,"\varepsilon_\rho"]\arrow[d,"\text{épi}"]&\{\pm1\}\arrow[r]\arrow[d,equal]&1\\
1\arrow[r]&\gamma'_{\rho,\sigma}\arrow[r]&\unsure{\gamma_{\rho,\sigma}|Z_{\rho,\sigma}}\arrow[r,"\varepsilon'_\gamma"]&\{\pm1\}\arrow[r]&1
\end{tikzcd}}
\]
où on a posé
\[
\tag{42}
\gamma'_{\rho,\sigma}=\operatorname{Ker}\bigl(
\diagram{\begin{tikzcd}
\gamma_{\rho,\sigma}\arrow[r,"\varepsilon'_\gamma"]&\{\pm1\}
\end{tikzcd}}\bigr);
\]

\deleted{b) Si $\sigma^2\ne1$ (cf (39)) alors on a commutativité}
\inserted{De $\overline{\mathrm m}$}
\[
\diagram{\deleted{\begin{tikzcd}[column sep=large,row sep=small]
N_\sigma^*\arrow[rr,"\delta_\sigma"]\arrow[dr,"\varepsilon_\sigma"']&&
\gamma_{\rho,\sigma}\arrow[dl,"\unsure{\varepsilon_\gamma|\varepsilon_\sigma}"]\\
&\{\pm1\}&
\end{tikzcd}}},
\]
\deleted{d'où} ce diagramme de suites exactes
\[
\tag{43}
\diagram{\begin{tikzcd}[column sep=large]
1\arrow[r]&Z_\sigma\arrow[r]\arrow[d,"\text{épi}"]&N_\sigma^*\arrow[r,"\varepsilon_\sigma"]\arrow[d,"\text{épi}"']&\{\pm1\}\arrow[r]\arrow[d,equal]&1\\
1\arrow[r]&\gamma''_{\rho,\sigma}\arrow[r]&\gamma_{\rho,\sigma}\arrow[r,"\unsure{\varepsilon_\gamma|\varepsilon_\sigma}"]&\{\pm1\}\arrow[r]&1
\end{tikzcd}},\qquad\text{où}
\]
\[
\tag{44}
\gamma''_{\rho,\sigma}=\operatorname{Ker}\bigl(
\diagram{\begin{tikzcd}
\gamma_{\rho,\sigma}\arrow[r,"\unsure{\varepsilon_\gamma|\varepsilon_\sigma}"]&\{\pm1\}
\end{tikzcd}}\bigr).
\]
\mspage{555}

{\small
\begin{center}
555\\
\texttt{https://grothendieck.umontpellier.fr}\\
Université de Montpellier, UMR CNRS 5149, Case courrier 051, Place Eugène Bataillon\\
34095 Montpellier cedex 5 - France
\end{center}
}

\scanpage{557}

\uline{NB} Les notations $\varepsilon_\gamma,\varepsilon'_\gamma:
\gamma_{\rho,\sigma}\longrightarrow\{\pm1\}$ peuvent prêter à confusion,
la première \unsure{référant|afférant} à l'action d'un $u$ sur $L_\sigma$
(non sur $L_\rho$). \unsure{Une} notation un peu plus lourde,
mais ne prêtant pas à confusion, serait
$\varepsilon_{\gamma,\rho},\varepsilon_{\gamma,\sigma}$ --- ou $\overline{\mathrm m}$
seulement $\varepsilon_\rho,\varepsilon_\sigma$, la confusion avec les hom de ce nom
$\unsure{N'_\rho|N_\rho^*}\longrightarrow\{\pm1\}$,
$N_\sigma\longrightarrow\{\pm1\}$ étant anodine lorsque
$\rho^2\ne1,\sigma^2\ne1$, à cause des compatibilités (41) (43).

On a donc construit, à l'aide \unsure{de|des} (1) i.e. de $\mathcal G$ muni de
$\rho,\sigma,\tau'$, des groupes \inserted{$\mathcal G^+$,}
$G_{\rho,\sigma},\gamma_{\rho,\sigma},Z_{\rho,\sigma},N'_\rho,N'_\sigma$,
s'insérant dans un diagramme commutatif de groupes
\[
\tag{45}
\diagram{\begin{tikzcd}[column sep=large,row sep=small]
&L_0\simeq SZ_{\rho,\sigma}\arrow[rr]\arrow[ddd,hook]&&Z_{\rho,\sigma}
\arrow[ddd,hook,dashed]\arrow[dddrr,dashed,"\delta_Z"]&&&\\
&&SN_\rho^*\arrow[rr]\arrow[ddl,hook]& &N_\rho^*
\arrow[ddl,hook,dashed]\arrow[ddr,dashed,"\delta_\rho"]&&
\widehat{\mathbb Z}^{*}\times\{\pm1\}\times\{\pm1\}\\
&&SN_\sigma^*\arrow[rr]\arrow[dl,hook,dotted]\arrow[dl]&&N_\sigma^*
\arrow[dl,hook,dotted]\arrow[dl]\arrow[dr,dashed,"\delta_\sigma"]&&\\
1\arrow[r]&\mathcal G^+\arrow[rr,hook,"i_{\mathcal G^+}"]\arrow[dr,hook]&&
G_{\rho,\sigma}\arrow[rr,"\delta_G"]&&\gamma_{\rho,\sigma}\arrow[r]\arrow[uur,"{(\chi_\gamma,\varepsilon_{\gamma,\rho},\varepsilon_{\gamma,\sigma})}"]&1\\
&&\mathcal G\arrow[ur,hook]&&&&
\end{tikzcd}}
\]
où on a posé
\[
\tag{46}
\left\{\begin{aligned}
SZ_{\rho,\sigma}&=i_{L_0,Z}(L_0)
=\operatorname{Ker}(\delta_Z:Z_{\rho,\sigma}\longrightarrow\gamma_{\rho,\sigma})\\
SN'_\rho&=\operatorname{Ker}(\delta_\rho:N_\rho^*\longrightarrow\gamma_{\rho,\sigma})
=\operatorname{Cent}_{\mathcal G^+}(\rho)\\
SN'_\sigma&=\operatorname{Ker}(\delta_\sigma:N_\sigma^*\longrightarrow\gamma_{\rho,\sigma})
=\operatorname{Cent}_{\mathcal G^+}(\sigma)
\end{aligned}\right.
\]
\margin{NB On peut donc aussi les noter $SZ_\rho,SZ_\sigma$ au lieu des
$SN_\rho^*,SN_\sigma^*$.}
\mspage{556}

{\small
\begin{center}
556\\
\texttt{https://grothendieck.umontpellier.fr}\\
Université de Montpellier, UMR CNRS 5149, Case courrier 051, Place Eugène Bataillon\\
34095 Montpellier cedex 5 - France
\end{center}
}

\scanpage{558}
\mspage{582}

\unsure{de} sorte qu'on a des suites exactes
\[
\tag{47}
\left\{
\diagram{\begin{tikzcd}[column sep=large]
1\arrow[r]&S\unsure{\deleted{Z}N}_\rho^*\arrow[r]&N_\rho^*\arrow[r,"\delta_\rho"]&
\gamma_{\rho,\sigma}\arrow[r]&1\\
1\arrow[r]&SN_\sigma^*\arrow[r]&N_\sigma^*\arrow[r,"\delta_\sigma"]&
\gamma_{\rho,\sigma}\arrow[r]&1
\end{tikzcd}}
\right.
\]
\margin{On a
\[
\tag{47bis}
\left\{\begin{aligned}
\operatorname{Im}\bigl(\unsure{SN_\rho}\longrightarrow\mathcal G^+\bigr)
&=\operatorname{Norm}_{\mathcal G^+}(L_\rho)\\
\operatorname{Im}\bigl(SN_\sigma\longrightarrow\mathcal G^+\bigr)
&=\operatorname{Norm}_{\mathcal G^+}(L_\sigma)
\end{aligned}\right.
\]
(ou encore, si $\rho^2\ne1$ resp. $\sigma^2\ne1$,
$N_\rho\cap\mathcal G$ et $N_\sigma\cap\mathcal G$)}
s'envoyant dans la suite exacte
\[
\diagram{\begin{tikzcd}[column sep=large]
1\arrow[r]&\mathcal G^+\arrow[r]&G_{\rho,\sigma}\arrow[r]&
\gamma_{\rho,\sigma}\arrow[r]&1
\end{tikzcd}},
\]
de façon analogue à (27). De plus, les hom
$i_\rho:N_\rho^*\longrightarrow G_{\rho,\sigma}$ et
$i_\sigma:N_\sigma^*\longrightarrow G_{\rho,\sigma}$ ont
une nette tendance à être injectifs (cf (39)),
ce qu'on a marqué sur le diagramme (45) par le signe
$\diagram{\begin{tikzcd}[column sep=1.7em]
\phantom{0}\arrow[r,hook,dotted]\arrow[r]&\phantom{0}
\end{tikzcd}}$
(signe d'inclusion pointillé).

On introduit aussi
\[
\tag{48}
\left\{\begin{aligned}
\chi_\rho&=\chi_\gamma\circ\delta_\rho:
\diagram{\begin{tikzcd}
\unsure{\deleted{Z}N}_\rho^*\arrow[r]&\widehat{\mathbb Z}^{*}
\end{tikzcd}}\\
\chi_\sigma&=\chi_\gamma\circ\unsure{\delta_{\deleted{\rho}\sigma}}:
\diagram{\begin{tikzcd}
\unsure{\deleted{Z}N}_\sigma^*\arrow[r]&\widehat{\mathbb Z}^{*}
\end{tikzcd}}
\end{aligned}\right.,
\]
dont les noyaux (qui contiennent $S\unsure{\deleted{Z}N}_\rho\deleted{\unsure{\sigma}}$
resp. $SN_\sigma$, et qu'on \unsure{se gardera} de confondre avec eux)
sont notés comme de juste $N_\rho^{*+},N_\sigma^{*+}$.

\uline{NB} Il faudrait encore examiner les composés
\[
\tag{49}
\left\{
\diagram{\begin{tikzcd}[column sep=large]
N_\rho^*\arrow[r]&\gamma_{\rho,\sigma}\arrow[r,"\unsure{\varepsilon_{\gamma,\sigma}}"]&\{\pm1\}\\
N_\sigma^*\arrow[r]&\gamma_{\rho,\sigma}\arrow[r,"\varepsilon_{\gamma,\rho}"]&\{\pm1\}
\end{tikzcd}}
\right.,
\]
on verra que dans les "bons cas", les restrictions de ces hom. à
$Z_\rho,Z_\sigma$ respectivement sont épi.
\mspage{557}

{\small
\begin{center}
557\\
\texttt{https://grothendieck.umontpellier.fr}\\
Université de Montpellier, UMR CNRS 5149, Case courrier 051, Place Eugène Bataillon\\
34095 Montpellier cedex 5 - France
\end{center}
}

\scanpage{559}
\mspage{\unsure{583}}

\deleted{On pose}
\[
\deleted{
\left\{
\begin{aligned}
S\unsure{Z_\rho|N_\rho}
 &=\operatorname{Ker}\bigl(\unsure{Z_\rho|N_\rho}
       \xrightarrow{\delta_\rho}\gamma_{\rho,\sigma}\bigr)
       =\operatorname{Norm}_{\mathcal G}(L_\rho)\\
SN_\sigma
 &=\operatorname{Ker}\bigl(N_\sigma
       \xrightarrow{\delta_\sigma}\gamma_{\rho,\sigma}\bigr)
       =\operatorname{Norm}_{\mathcal G}(L_\sigma)
\end{aligned}
\right.
}
\tag{\unsure{15}}
\]
\[
\deleted{
S\unsure{B}_{\rho,\sigma}
=\operatorname{Ker}\bigl(\unsure{B}_{\rho,\sigma}
   \xrightarrow{\delta_B}\gamma_{\rho,\sigma}\bigr)
\simeq L_0\subset
\left\{
\begin{gathered}
\unsure{\operatorname{Norm}|\operatorname{Cent}}_{\mathcal G}
       (\unsure{L_0|L_\rho})\\[-2pt]
\cap\\[-2pt]
\operatorname{Norm}_{\mathcal G}\,\unsure{L_0|L_\sigma}
\end{gathered}
\right.
}
\]
\margin{Annulé\par
\deleted{\unsure{En réalité que}
\unsure{dans le cas que} \unsure{on veut considérer},
\(\operatorname{Norm}_{\mathcal G}(L_0)
\unsure{\ne|=}\unsure{\operatorname{Cent}}_{\mathcal G}(L_0)\)
et idem pour \(L_\rho,L_\sigma\),}}

\textcircled{\(5^\circ\)}
\margin{Les groupes
\(S_0\mathcal M_{\rho,\sigma},\mathcal M_{\rho,\sigma},
S_0\Gamma_{\rho,\sigma},\Gamma_{\rho,\sigma}\).}

On pose
\[
\left\{
\begin{aligned}
S_0\mathcal M_{\rho,\sigma}
 &=Z_{\rho,\sigma}
   \mathop{\times}_{\gamma_{\rho,\sigma}}N_\rho^*
   \mathop{\times}_{\gamma_{\rho,\sigma}}N_\sigma^*
   \mathop{\times}_{\gamma_{\rho,\sigma}}G_{\rho,\sigma}.\\
\mathcal M_{\rho,\sigma}
 &=N_\rho^*
   \mathop{\times}_{\gamma_{\rho,\unsure{\sigma}}}N_\sigma^*
   \mathop{\times}_{\gamma_{\rho,\sigma}}G_{\rho,\sigma}
\end{aligned}
\right.
\tag{50}
\]
\[
\deleted{\unsure{42}}
\qquad
\left\{
\begin{aligned}
S_0\Gamma_{\rho,\sigma}
 &=Z_{\rho,\sigma}
   \mathop{\times}_{\gamma_{\rho,\sigma}}N_\rho^*
   \mathop{\times}_{\gamma_{\rho,\sigma}}N_\sigma^*\\
\Gamma_{\rho,\sigma}
 &=\unsure{N_\rho^*}
   \mathop{\times}_{\gamma_{\rho,\sigma}}N_\sigma^*,
\end{aligned}
\right.
\tag{51}
\]

d'où \deleted{\unsure{hom.}}\inserted{\uline{diagr.}} de suites exactes
\diagram{
\((52)\qquad\)
\begin{tikzcd}[column sep=large,row sep=large]
 & & 1 \arrow[d] & 1 \arrow[d] & \\
 & & L_0 \arrow[d] \arrow[r,equal] & L_0 \arrow[d] & \\
1 \arrow[r] & \mathcal G^+ \arrow[r] \arrow[d,equal]
 & S_0\mathcal M_{\rho,\sigma} \arrow[r] \arrow[d]
 & S_0\Gamma_{\rho,\sigma} \arrow[r] \arrow[d] & 1\\
1 \arrow[r] & \mathcal G^+ \arrow[r]
 & \mathcal M_{\rho,\sigma} \arrow[r] \arrow[d]
 & \Gamma_{\rho,\sigma} \arrow[r] \arrow[d] & 1,\\
 & & 1 & 1 &
\end{tikzcd}
}

et des suites exactes
\diagram{
\((\unsure{53|54})\qquad\)
\begin{tikzcd}[column sep=large]
1 \arrow[r]
 & \underset{\simeq L_0}{\uline{SZ_{\rho,\sigma}}}
   \times SN_\rho^*\times SN_\sigma^*\times\mathcal G^+
       \arrow[r]
 & S_0\mathcal M_{\rho,\sigma}\arrow[r]
 & \gamma_{\rho,\sigma}\arrow[r] & 1
\end{tikzcd}
}
\diagram{
\((54)\qquad\)
\begin{tikzcd}[column sep=large]
1 \arrow[r]
 & SN_\rho^*\times SN_\sigma^*\times\mathcal G^+\arrow[r]
 & \mathcal M_{\rho,\sigma}\arrow[r]
 & \gamma_{\rho,\sigma}\arrow[r] & 1
\end{tikzcd}
}
\diagram{
\((55)\qquad\)
\begin{tikzcd}[column sep=large]
1 \arrow[r]
 & \underset{\simeq L_0}{\uline{SZ_{\rho,\sigma}}}
   \times SN_\rho^*\times SN_\sigma^*\arrow[r]
 & S_0\Gamma_{\rho,\sigma}\arrow[r]
 & \gamma_{\rho,\sigma}\arrow[r] & 1
\end{tikzcd}
}
\diagram{
\((56)\qquad\)
\begin{tikzcd}[column sep=large]
1 \arrow[r]
 & SN_\rho^*\times SN_\sigma^*\arrow[r]
 & \Gamma_{\rho,\sigma}\arrow[r]
 & \gamma_{\rho,\sigma}\arrow[r] & 1.
\end{tikzcd}
}

\mspage{558}

{\small
\begin{center}
558\\
\texttt{https://grothendieck.umontpellier.fr}\\
Université de Montpellier, UMR CNRS 5149, Case courrier 051, Place Eugène Bataillon\\
34095 Montpellier cedex 5 - France
\end{center}
}

\scanpage{560}

On a aussi
\[
S_0\mathcal M_{\rho,\sigma}
\simeq\mathcal M_{\rho,\sigma}
       \mathop{\times}_{\gamma_{\rho,\sigma}}Z_{\rho,\sigma}
\tag{57}
\]
\[
S_0\Gamma_{\rho,\sigma}
\simeq\Gamma_{\rho,\sigma}
       \mathop{\times}_{\gamma_{\rho,\sigma}}Z_{\rho,\sigma}
\tag{58}
\]

i.e. les ext. \(S_0\mathcal M_{\rho,\sigma},S_0\Gamma_{\rho,\sigma}\)
de \(\mathcal M_{\rho,\sigma},\Gamma_{\rho,\sigma}\) par \(L_0\)
se déduisent de l'extension \(Z_{\rho,\sigma}\) de
\(\gamma_{\rho,\sigma}\) par \(L_0\) par image inverse par
\(\mathcal M_{\rho,\sigma}\to\gamma_{\rho,\sigma},
\Gamma_{\rho,\sigma}\to\gamma_{\rho,\sigma}\) --- et on a aussi
\deleted{\unsure{pour}}
\[
S_0\mathcal M_{\rho,\sigma}
\simeq\mathcal M_{\rho,\sigma}
       \mathop{\times}_{\Gamma_{\rho,\sigma}}S_0\Gamma_{\rho,\sigma}
\tag{59}
\]

i.e. l'extension \(S_0\mathcal M_{\rho,\sigma}\) de
\(\mathcal M_{\rho,\sigma}\) par \(L_0\) se déduit de l'extension
\(S_0\Gamma_{\rho,\sigma}\) de \(\Gamma_{\rho,\sigma}\) par \(L_0\),
par image inverse \unsure{par}
\(\mathcal M_{\rho,\sigma}\to\Gamma_{\rho,\sigma}\).

\inserted{On introduit un élément}
\[
\deleted{\unsure{\tau'_{\mathcal M}}},
\qquad
\tau'_{\deleted{\unsure{S_0}}\mathcal M}
  \in\deleted{S_0}\mathcal M_{\rho,\sigma}
\quad
\deleted{\text{\unsure{les}}}
\quad
\tau'_{\deleted{\unsure{S_0}}\mathcal M}
 =\deleted{\unsure{\tau'_Z}}
 \bigl(
  \overset{\inserted{\substack{a\\\Vert}}}{\tau'_\rho},
  \overset{\inserted{\substack{b\\\Vert}}}{\tau'_\sigma},
  \overset{\inserted{\substack{u\\\Vert}}}{\tau'_G}
 \bigr)
\tag{60}
\]

\deleted{\unsure{d'homomorphismes} entre les groupes
introduits, je vais les relier dans
\deleted{dans}\inserted{un} diagramme commutatif}
\inserted{où \(\deleted{\unsure{\tau'_Z}},\tau'_\rho,\tau'_\sigma,
\tau'_G\) sont les}
"avatars" de \(\tau'\), pris resp. dans
\(\deleted{\unsure{Z_{\rho,\sigma}}}
N_\rho^*,N_\sigma^*,G_{\rho,\sigma}\).
Ceux-ci ont bien tous la \(\overline{\mathrm m}\) image dans
\(G_{\rho,\sigma}\) et a fortiori dans \(\gamma_{\rho,\sigma}\),
donc définissent un \unsure{élément} de
\(\deleted{\unsure{G}}\inserted{\mathcal M}_{\rho,\sigma}\),
dont l'opération sur \(\mathcal G^+\) est celle de \(\tau\).
D'où un plongement
\diagram{
\((61)\qquad\)
\begin{tikzcd}
\mathcal G \arrow[r,hook,"i_{\mathcal G,\mathcal M}"]
 & \mathcal M_{\rho,\sigma}
\end{tikzcd}
\quad : identité sur \(\mathcal G^+\)
\qquad
\begin{tikzcd}
\tau'_{\mathcal G}\arrow[r,mapsto] & \tau'_{\mathcal M}.
\end{tikzcd}
}

\mspage{559}

{\small
\begin{center}
559\\
\texttt{https://grothendieck.umontpellier.fr}\\
Université de Montpellier, UMR CNRS 5149, Case courrier 051, Place Eugène Bataillon\\
34095 Montpellier cedex 5 - France
\end{center}
}

\scanpage{561}
\mspage{584}

\diagram{
\((62)\qquad\)
\begin{tikzcd}[column sep=small,row sep=large]
 & & L_0\times SN_\rho^*\times SN_\sigma^*
 & & & & L_0 & & & &\\
 & & & & S_0\Gamma_{\rho,\sigma}
 & & & & Z_{\rho,\sigma} & &\\
L_0\times SN_\rho^*\times SN_\sigma^*\times\mathcal G^+
 & & SN_\rho^*\times SN_\sigma^*
 & & & & \mathrm{III}^{\unsure{+|'}}
 & & & &\\
 & & & S_0\mathcal M_{\rho,\sigma}
 & \Gamma_{\rho,\sigma} & & N_\rho^*
 & & & & \{\pm1\}\\
 & SN_\rho^*\times SN_\sigma^*\times\mathcal G^+
 & & & & & N_\sigma^*
 & & \gamma_{\rho,\sigma} & & \widehat{\mathbb Z}^{*}\\
 & & & & & \mathrm{II}
 & \deleted{\unsure{\mathcal G^+|N_\sigma^*}}
 & & & & \{\pm1\}\\
 & & & \mathcal M_{\rho,\sigma}
 & & & & & G_{\rho,\sigma} & &\\
\mathcal G^+ & & & & & & & & & &\\
 & & \mathcal G & & & & & & & &\\
 & & \Vert & & & & & & & &\\
 & & \{1,\tau'\}\mathbin{.}\mathcal G^+
 & & & & & & & &
\arrow[from=1-3,to=1-7,"\mathrm{pr}"]
\arrow[from=1-3,to=2-5,hook]
\arrow[from=1-3,to=3-3,"\mathrm{pr}"]
\arrow[from=1-7,to=2-9,hook]
\arrow[from=2-5,to=2-9,"\theta_Z"]
\arrow[from=2-5,to=4-5]
\arrow[from=2-9,to=5-9,"\delta_Z"]
\arrow[from=3-1,to=1-3,"\mathrm{pr}"]
\arrow[from=3-1,to=5-2,"\mathrm{pr}"']
\arrow[from=3-1,to=4-4,hook]
\arrow[from=3-3,to=4-5,hook]
\arrow[from=4-4,to=2-5]
\arrow[from=4-4,to=7-4]
\arrow[from=4-5,to=4-7,"\theta_\rho"]
\arrow[from=4-5,to=5-7,"\theta_\sigma"']
\arrow[from=4-7,to=5-9,"\delta_\rho"]
\arrow[from=5-7,to=5-9,"\delta_\sigma"]
\arrow[from=5-2,to=3-3,"\mathrm{pr}"']
\arrow[from=5-2,to=7-4,hook]
\arrow[from=5-2,to=6-7,dashed,"\unsure{\mathrm{pr}}"]
\arrow[from=6-7,to=7-9,dashed,hook]
\arrow[from=7-4,to=4-5]
\arrow[from=7-4,to=6-7,dotted]
\arrow[from=7-4,to=7-9,"\theta_G"]
\arrow[from=7-9,to=5-9,"\delta_G"']
\arrow[from=5-9,to=4-11,
 "{\varepsilon_{\gamma,\rho}=\varepsilon'_\gamma}"]
\arrow[from=5-9,to=5-11,"\chi_\gamma"]
\arrow[from=5-9,to=6-11,
 "{\varepsilon_{\gamma,\sigma}=\varepsilon_\gamma}"']
\arrow[from=8-1,to=5-2,hook,"\uline{\mathrm{inj}}"]
\arrow[from=8-1,to=7-4,dashed]
\arrow[from=8-1,to=9-3,hook]
\arrow[from=9-3,to=7-4,hook,"i_{\mathcal G,\mathcal M}"]
\arrow[from=9-3,to=7-9,dashed,hook,"i_{\mathcal G,G}"']
\arrow[from=1-3,to=2-9,phantom,
 "{\deleted{\unsure{\mathrm{IV}}}\ \mathrm{V}}" description,pos=.52]
\arrow[from=4-7,to=5-7,phantom,"\mathrm{I}" description]
\arrow[from=8-1,to=7-4,phantom,"\mathrm{IV}" description,pos=.5]
\end{tikzcd}
}

On a résumé \deleted{\unsure{les prin.}}\inserted{certaines des}
relations entre les groupes introduits jusqu'à présent dans le diagramme
commutatif (62), que je commente brièvement.

a) Toutes les flèches sont soit injections \((\hookrightarrow)\)
soit surjections \((\longrightarrow)\), à l'exception
\inserted{peut-être} des trois flèches issues de \(\gamma_{\rho,\sigma}\).

b) Les dix carrés (six faces d'un cube, et
\deleted{quatre}\inserted{\unsure{trois|les}} carrés marqués
\(\mathrm I,\mathrm{II},\mathrm{III}^{\unsure{+|'}},
\mathrm{IV},\mathrm V\)) sont cartésiens.

c) Les \inserted{\unsure{cinq}} flèches marquées \(\mathrm{pr}\)
sont des projections, canoniques de groupes produits sur des produits
partiels. La flèche marquée \(\mathrm{inj}\)
\diagram{
\begin{tikzcd}
\mathcal G^+\arrow[r] &
SN_\rho^*\times SN_\sigma^*\times\mathcal G^+
\end{tikzcd}
}
est l'injection canonique.

d) Pour ne pas détruire la commutativité du diagramme,
nous avons omis d'y faire figurer les flèches
\inserted{(d'inclusion, notamment)}
\diagram{
\begin{tikzcd}
Z_{\rho,\sigma}\arrow[r] & G_{\rho,\sigma},
\end{tikzcd}
\qquad
\begin{tikzcd}
N_\rho^*\arrow[r] & G_{\rho,\sigma},
\end{tikzcd}
\qquad
\begin{tikzcd}
N_\sigma^*\arrow[r] & G_{\rho,\sigma}
\end{tikzcd}
}
qui figuraient dans (45).

\mspage{560}

{\small
\begin{center}
560\\
\texttt{https://grothendieck.umontpellier.fr}\\
Université de Montpellier, UMR CNRS 5149, Case courrier 051, Place Eugène Bataillon\\
34095 Montpellier cedex 5 - France
\end{center}
}

\scanpage{562}

e) Pour ne pas surcharger le diagramme, je n'ai pas introduit dans
le diagramme des notations pour certaines flèches composées utiles,
telles
\[
(63)\quad
\left\{
\begin{array}{ll}
\delta_\Gamma:\Gamma_{\rho,\sigma}\longrightarrow\gamma_{\rho,\sigma},
&
\deleted{\unsure{\delta_{S_0\Gamma}}:
\unsure{S_0\Gamma_{\rho,\sigma}}\longrightarrow
\unsure{\gamma_{\rho,\sigma}}.}\\
\delta_{S_0\Gamma}:S_0\Gamma_{\rho,\sigma}
\longrightarrow\gamma_{\rho,\sigma}&\\
\delta_{\mathcal M}:\mathcal M_{\rho,\sigma}
\longrightarrow\gamma_{\rho,\sigma}&\\
\delta_{S_0\mathcal M}:S_0\mathcal M_{\rho,\sigma}
\longrightarrow\gamma_{\rho,\sigma}&
\end{array}
\right.
\]
et
\[
(64)\quad
\left\{
\begin{aligned}
\chi_\Gamma&=\chi_\gamma\circ\delta_\Gamma:
\Gamma_{\rho,\sigma}\longrightarrow\widehat{\mathbb Z}^{*}\\
\chi_{S_0\Gamma}&=\chi_\gamma\circ\delta_{S_0\Gamma}:
S_0\Gamma_{\rho,\sigma}\longrightarrow\widehat{\mathbb Z}^{*}\\
\chi_{\mathcal M}&=\chi_\gamma\circ\delta_{\mathcal M}:
\mathcal M_{\rho,\sigma}\longrightarrow\widehat{\mathbb Z}^{*}\\
\chi_{S_0\mathcal M}&=\chi_\gamma\circ\delta_{S_0\mathcal M}:
S_0\mathcal M_{\rho,\sigma}\longrightarrow\widehat{\mathbb Z}^{*}
\end{aligned}
\right.
\]
et les composés
\(\varepsilon_{?,\rho},\varepsilon_{?,\sigma}:?\longrightarrow\{\pm1\}\),
où \(?\) peut prendre \deleted{\unsure{toute}} valeur
\(Z,G\) (déjà introduits) \(N_\rho^*,N_\sigma^*\)
(on aura
\(\varepsilon_{N_\rho,\rho}=\varepsilon_\rho
\deleted{\text{\unsure{simplement}}},\quad
\varepsilon_{N_\sigma,\sigma}=\varepsilon_\sigma
\deleted{\text{\unsure{simplement}}}\)),
\(\Gamma,S_0\Gamma,\mathcal M,S_0\mathcal M\).
Le noyau de \(\chi_?\) est noté \(?^+\), celui de
\(\varepsilon_{?,\rho}\) est noté \(?'\), celui de
\(\varepsilon_{?,\sigma}\) est noté \(?''\),
on note
\[
?^\circ\doteq?'\cap?'',\qquad
?^{+'}=?^+\cap?',\quad
?^{+''}=?^+\cap?'',
\]
\[
?^{+,\circ}=?^+\cap?^\circ,\quad
S?=\operatorname{Ker}\delta_?,
\]
\deleted{et} donc \(\deleted{\unsure{S?\subset}}\)

\margin{\deleted{il y a une petite incompatibilité avec la notation
\(N'_\rho,N'_\sigma\) déjà utilisée et \unsure{devenue} entre temps
\unsure{fausse}.}}

\[
(65)
\]
\diagram{
\begin{tikzcd}[row sep=1.15em,column sep=1.6em]
 & & ?^{'+} & &\\
\deleted{\unsure{S?\subset}}\quad ?^{\circ+}
 & & & & ?^+\ \deleted{\text{\unsure{on}}}\\
 & & ?^{''+} & &\\
 & & \unsure{?'} & &\\
?^\circ & & & & ?\\
 & & ?'' & &
\arrow[from=2-1,to=1-3,hook]
\arrow[from=1-3,to=2-5,hook]
\arrow[from=2-1,to=3-3,hook]
\arrow[from=3-3,to=2-5,hook]
\arrow[from=2-1,to=5-1,hook]
\arrow[from=1-3,to=4-3,hook,shift left=.7ex]
\arrow[from=3-3,to=6-3,hook,shift right=.7ex]
\arrow[from=2-5,to=5-5,hook]
\arrow[from=5-1,to=4-3,hook]
\arrow[from=5-1,to=6-3,hook]
\arrow[from=4-3,to=5-5,hook]
\arrow[from=6-3,to=5-5,hook]
\end{tikzcd}
}
cube "cartésien".

f) Rappelons aussi qu'on a commutativité dans
\[
(66)
\]
\diagram{
\begin{tikzcd}[row sep=.8em,column sep=2em]
 & \{\pm1\}\arrow[dr] &\\
\mathcal G
\arrow[ur,"{\varepsilon_{\mathcal G,\rho}}"]
\arrow[rr,"\chi_{\mathcal G}"]
\arrow[dr,"{\varepsilon_{\mathcal G,\sigma}}"']
 & & \widehat{\mathbb Z}^{*}\\
 & \{\pm1\}\arrow[ur] &
\end{tikzcd}
}
, soit \(\varepsilon_{\mathcal G}:\mathcal G\longrightarrow\{\pm1\}\),

\mspage{561}

{\small
\begin{center}
561\\
\texttt{https://grothendieck.umontpellier.fr}\\
Université de Montpellier, UMR CNRS 5149, Case courrier 051, Place Eugène Bataillon\\
34095 Montpellier cedex 5 - France
\end{center}
}

\scanpage{563}
\mspage{585}

Le noyau de $\varepsilon_{\unsure{\mathcal{G}|\mathfrak{G}}}$ n'est autre que $\mathcal{G}^{+}$. En \unsure{tant} que sous-groupe de $\mathcal M_{\rho,\sigma}$ (resp. de $G_{\rho,\sigma}$) $\mathcal{G}$ s'identifie à l'image inverse dans ce groupe du sous-groupe $\{1,\tau'_{\Gamma}\}$ de $\Gamma_{\rho,\sigma}$ (resp. \inserted{du ss-groupe} $\{1,\tau'_{\gamma}\}$ de $\gamma_{\rho,\sigma}$), où $\tau'_{\Gamma},\tau'_{\gamma}$ désignent les images de $\tau'\in\mathcal{G}$ dans $\Gamma_{\rho,\sigma}$ et $\gamma_{\rho,\sigma}$ respectivement.

g) Le diagramme (62) a comme partie principale
\[
\text{(67)}\qquad
\diagram{\begin{tikzcd}[row sep=0.9em,column sep=1.25em]
{} & {}\arrow[rrrr]\arrow[ddr] & {} & {} & {} & {}\arrow[dd] & {} & {} & {} & {}\\
{} & {} & {} & {}\arrow[drr] & {} & {} & {} & {} & {} & {}\\
{}\arrow[uur]\arrow[ddr] & {} & {}\arrow[ur]\arrow[dr] & {} & {} & {}\arrow[urrr,dashed,shorten <=3em]\arrow[rrr,shorten <=3em]\arrow[drrr,dashed,shorten <=3em] & \Biggl[ & {} & {} & \Biggr]\\
{} & {} & {} & {}\arrow[urr] & {} & {} & {} & {} & {} & {}\\
\deleted{\unsure{[\mathcal{G}^{+}]}} & {}\arrow[uur]\arrow[rrrr] & {} & {} & {} & {}\arrow[uu] & {} & {} & {} & {}\\
\unsure{\Bigl[}{}\arrow[ur,dashed,hook]\arrow[urrrrr,dashed,hook] & {} & {} & {} & {} & \Bigr] & {} & {} & {} & {}
\end{tikzcd}},
\]
tous les groupes marqués sur ce \inserted{dernier} diagramme, à l'exception de $\widehat{\mathbb Z}^{*},\{\pm1\},\{\pm1\}$ (mis entre crochets $[\;]$) s'envoient dans $\gamma_{\rho,\sigma}$, qui est le terme terminal de ce diagramme, reliant \unsure{9} groupes (en ne \unsure{comptant} pas $\mathcal{G}\ \deleted{\unsure{\mathcal{G}^{+}}}$ insérés \deleted{\unsure{par}} comme \deleted{\unsure{termes}} \deleted{\unsure{initiaux}} \inserted{en dehors} (flèches pointillées) dont $\gamma_{\rho,\sigma}$, déduits de quatre d'entre eux $Z_{\rho,\sigma},N_{\rho}^{*},N_{\sigma}^{*},G_{\rho,\sigma}$ (qui s'envoient dans $\gamma_{\rho,\sigma}$ par des épi) en prenant des produits fibrés sur $\gamma_{\rho,\sigma}$ de certains d'entre eux --- \deleted{\unsure{le}} le terme initial étant $S_{0}\mathcal M_{\rho,\sigma}$, qui est le produit fibré des quatre facteurs précédents. En prenant les noyaux des hom. de ces \deleted{quatre} neuf groupes dans $\gamma_{\rho,\sigma}$, on trouve un diagramme de $\overline{\mathrm m}$ type combinatoire, reliant

\mspage{562}

{\small
\begin{center}
562\\
\texttt{https://grothendieck.umontpellier.fr}\\
Université de Montpellier, UMR CNRS 5149, Case courrier 051, Place Eugène Bataillon\\
34095 Montpellier cedex 5 - France
\end{center}
}

\scanpage{564}

neuf groupes, \deleted{dont $\{1\}=\operatorname{Ker}(\gamma_{\rho,\sigma}\longrightarrow\gamma_{\rho,\sigma})$,} formés de $L_{0}=SZ_{\rho,\sigma}$, $SN_{\rho}^{*},SN_{\sigma}^{*}$ et $\mathcal{G}^{+}=SG_{\rho,\sigma}$ et de certains produits de ces groupes (dont le produit sur l'ens d'indices vide, savoir $\{1\}=\operatorname{Ker}(\gamma_{\rho,\sigma}\xrightarrow{\mathrm{id}}\gamma_{\rho,\sigma})$), avec comme \unsure{flèchismes|flèches hom.} des projections canoniques. Nous avons inséré dans (62) seulement \deleted{cinq} parmi ces groupes produits, avec les injections canoniques dans les produits fibrés dont ils proviennent, \unsure{en nous} dispensant d'y faire figurer séparément $SN_{\rho}^{*},SN_{\sigma}^{*},SG_{\rho,\sigma}=\mathcal{G}^{+}$, et les inclusions canoniques $SN_{\rho}^{*}\hookrightarrow N_{\rho}^{*}$, $SN_{\sigma}^{*}\hookrightarrow N_{\sigma}^{*}$, $\mathcal{G}^{+}=SG_{\rho,\sigma}\hookrightarrow G_{\rho,\sigma}$, et le \deleted{\unsure{p}} diagramme \inserted{complet} des projections des produits partiels de ces groupes et de $L_{0}$, \deleted{y compris} \inserted{où \unsure{manquent} donc} les projections sur ces \inserted{trois} groupes, et les hom. de ces groupes et de $L_{0}$ dans le groupe unité.

\mspage{563}

{\small
\begin{center}
563\\
\texttt{https://grothendieck.umontpellier.fr}\\
Université de Montpellier, UMR CNRS 5149, Case courrier 051, Place Eugène Bataillon\\
34095 Montpellier cedex 5 - France
\end{center}
}

\scanpage{565}
\mspage{586}

6°) \uline{L'ensemble} $\Sigma_{\rho,\sigma},S\unsure{\mathcal{G}|g}_{\rho,\sigma}$.

On pose
\[
\text{(68)}\qquad
\begin{gathered}
\Sigma_{\rho,\sigma}\unsure{\subset|=}\mathcal{G}^{+}\times\mathcal{G}^{+}\times\widehat{\mathbb Z}^{*}\times\{\pm1\}\times\{\pm1\}\\
\begin{array}{c}
\Vert\\[-0.4em]
\displaystyle
\left\{(\xi,\eta,\mu,\unsure{\varepsilon',\varepsilon})\;\middle|\;
\begin{array}{l}
\xi\rho\ \inserted{\mathcal{G}^{+}\text{-}}\text{conjugué à }\rho^{\varepsilon'}\\
\eta\sigma\ \mathcal{G}^{+}\text{-conjugué à }\sigma^{\varepsilon}\\
\xi=\eta\underbrace{\varepsilon_{1}^{\mu-1}}_{\rho_{1}^{\nu}\ \text{où }\nu=\frac{\mu-1}{2}}\quad\deleted{\unsure{\text{(avec}}}
\end{array}\right\}
\end{array}
\end{gathered}
\]

\uline{NB} Si $\rho$ non conjugué à $\rho^{-1}$ dans $\mathcal{G}^{+}$, $\varepsilon'$ connu quand on connaît $\xi,\eta$ --- de $\overline{\mathrm m}$ si $\sigma$ non conjugué à $\sigma^{-1}$ (donc $\sigma^{2}\ne1$), $\unsure{\varepsilon'|\varepsilon}$ est déterminé par $\xi,\eta$. \unsure{Si ce h.} $\widehat{\mathbb Z}\longrightarrow L_{0}$ $(n\longmapsto\varepsilon_{0}^{n})$ est injectif, alors $\mu$ est déterminé par $\xi,\eta$, et la relation $\xi=\eta\varepsilon_{1}^{\mu-1}$ s'écrit aussi (sans faire intervenir $\mu$)
\[
\text{(69)}\qquad \eta^{-1}\xi\in L_{1}.
\]

Soit d'autre part
\[
\text{(70)}\qquad
\begin{gathered}
S\mathcal{G}_{\rho,\sigma}\subset\mathcal{G}\times\mathcal{G}\times\widehat{\mathbb Z}^{*}\\
\begin{array}{c}
\Vert\\[-0.4em]
\displaystyle\left\{(\alpha,\beta,\mu)\;\middle|\;
\underbrace{[\alpha,\rho]}_{\xi}=[\beta,\sigma]\varepsilon_{1}^{\mu-1}\right\}
\end{array}
\end{gathered}
\]

On a une application
\[
\text{(71)}\qquad
\left\{\diagram{\begin{tikzcd}[row sep=1.5em,column sep=3em]
S\mathcal{G}_{\rho,\sigma}\arrow[r] & \Sigma_{\rho,\sigma}\\
(\alpha,\beta,\mu)\arrow[r,mapsto] & (\xi,\eta,\mu,\varepsilon',\varepsilon)
\end{tikzcd}}\right.
\quad\text{avec}\quad
\left\{\begin{array}{l}
\xi=[\alpha,\rho]\\
\eta=[\beta,\sigma]\\
\mu\ \text{le même}\\
\varepsilon'=\unsure{\varepsilon}_{\mathcal{G}}(\alpha)\\
\varepsilon=\varepsilon_{\mathcal{G}}(\beta)
\end{array}\right.
\]

\margin{\uline{NB} $\xi=\alpha\rho(\alpha)^{-1}\in\mathcal{G}^{+}$, car on a $\varepsilon_{\mathcal{G}}(\xi)=\unsure{\dot{\alpha}|\bar{\alpha}}\rho(\unsure{\dot{\alpha}|\bar{\alpha}})=1$ car $\rho$ opère trivialement sur $\mathcal{G}/\mathcal{G}^{+}\simeq\{\pm1\}$, pour la $\overline{\mathrm m}$ raison $\eta=\beta\sigma(\beta)^{\unsure{-1}}\in\mathcal{G}^{+}$.}

qui est \inserted{évidemment} \uline{surjective}.

On définit d'autre part

\mspage{564}

{\small
\begin{center}
564\\
\texttt{https://grothendieck.umontpellier.fr}\\
Université de Montpellier, UMR CNRS 5149, Case courrier 051, Place Eugène Bataillon\\
34095 Montpellier cedex 5 - France
\end{center}
}

\scanpage{566}

\[
\text{(72)}\qquad
\begin{gathered}
\Sigma_{\rho,\sigma}^{\mathrm{eff}}\subset\Sigma_{\rho,\sigma}\\
\begin{array}{c}
\Vert\\[-0.4em]
\displaystyle\left\{(\xi,\unsure{\eta},\mu,\varepsilon',\varepsilon)\in\Sigma_{\rho,\sigma}\;\middle|\;
\begin{array}{l}
\exists u\in\operatorname{Aut}(\mathcal{G}^{+})\ \text{avec}\\
u(\rho)=\xi\rho,\quad u(\unsure{\sigma})=\eta\sigma
\end{array}\right\}
\end{array}
\end{gathered}
\]
\margin{i.e. $[u,\rho]=\xi$, $[u,\sigma]=\eta$}
(cet $u$ sera alors unique, \inserted{et \unsure{inversement}, la \unsure{conn.} de $u$ détermine $\xi,\eta$.}) Il y a correspondance $1$-$1$ entre les élts de $\Sigma_{\rho,\sigma}^{\mathrm{eff}}$, et les systèmes $(u,\mu,\varepsilon',\varepsilon)\in\operatorname{Aut}(\mathcal{G}^{+})\times\widehat{\mathbb Z}^{*}\times\{\pm1\}\times\{\pm1\}$, tels que l'on ait \deleted{$u(\rho)\ \mathcal{G}^{+}\text{-conjug}\unsure{\text{.}}$}
\[
\left\{\begin{array}{lll}
u(\rho)\text{ est}&\mathcal{G}^{+}\text{-conjugué à}&\rho^{\varepsilon'}\\
u(\sigma)&\rule{9em}{0.35pt}&\sigma^{\varepsilon}\\
u(\varepsilon_{0})=\varepsilon_{0}^{\mu}&\multicolumn{2}{l}{(\text{qui exprime }\xi=\eta\rho_{1}^{\nu})}
\end{array}\right.
\]
--- or cet ensemble n'est autre que le sous-groupe $\unsure{Z_{\rho,\sigma}|\mathbb Z_{\rho,\sigma}}$ de $\operatorname{Aut}(\mathcal{G}^{+})\times\widehat{\mathbb Z}^{*}\times\{\pm1\}\times\{\pm1\}$. \inserted{(cf (6) p.\unsure{577'})}

On trouve ainsi une bijection
\[
\text{(73)}\qquad
\diagram{\begin{tikzcd}[row sep=2em,column sep=3em]
Z_{\rho,\sigma}\arrow[r,"\sim"] & \Sigma_{\rho,\sigma}^{\mathrm{eff}}\\
\begin{array}{c}\deleted{\unsure{(u,\mu,\varepsilon,\varepsilon')}}\\u\end{array}\arrow[r,mapsto] & (\xi,\eta,\mu,\varepsilon',\varepsilon)
\end{tikzcd}}
\quad\left\{\begin{array}{l}
\xi=[u,\rho]\\
\eta=[u,\sigma]\\
\mu=\chi(u)\\
\varepsilon'=\varepsilon'_{Z}(u)\\
\varepsilon=\varepsilon_{Z}(u)
\end{array}\right.\ .
\]

On désigne par $S\mathcal{G}_{\rho,\sigma}^{\mathrm{eff}}$ l'image inverse de $\Sigma_{\rho,\sigma}^{\mathrm{eff}}$ dans $S\mathcal{G}_{\rho,\sigma}$ par (71), donc on aura
\[
\text{(74)}\qquad
\begin{gathered}
S\mathcal{G}_{\rho,\sigma}^{\mathrm{eff}}\subset S\mathcal{G}_{\rho,\sigma}\\
\begin{array}{c}
\Vert\\[-0.4em]
\displaystyle\left\{\unsure{\{\alpha,\beta,\mu\}}\in\mathcal{G}\times\mathcal{G}\times\widehat{\mathbb Z}^{*}\;\middle|\;
\begin{array}{l}
[\alpha,\rho]=[\beta,\sigma]\varepsilon_{1}^{\mu-1}\\
\exists u\in\operatorname{Aut}(\mathcal{G}^{+})\ \text{tel qu'on ait}\\
\left\{\begin{array}{ll}
u(\rho)=\alpha(\rho)&(=\xi\rho,\text{ où }\xi=[\alpha,\rho])\\
u(\sigma)=\beta(\sigma)&(=\eta\sigma,\text{ où }\eta=[\beta,\sigma])
\end{array}\right.
\end{array}\right\}
\end{array}
\end{gathered}
\]
de sorte qu'on a un \unsure{hom.} \uline{surjectif} induit par (71)
\[
\text{(75)}\qquad
\left|\diagram{\begin{tikzcd}[column sep=3em]
S\mathcal{G}_{\rho,\sigma}^{\mathrm{eff}}\arrow[r]&\Sigma_{\rho,\sigma}^{\mathrm{eff}}
\end{tikzcd}}\right.\ .
\]

\mspage{565}

{\small
\begin{center}
565\\
\texttt{https://grothendieck.umontpellier.fr}\\
Université de Montpellier, UMR CNRS 5149, Case courrier 051, Place Eugène Bataillon\\
34095 Montpellier cedex 5 - France
\end{center}
}

\scanpage{567}
\mspage{587}

On a une application canonique
\[
\text{(76)}\qquad
\left\{\diagram{\begin{tikzcd}[row sep=2em,column sep=2em]
S_{0}\Gamma_{\rho,\sigma}=\unsure{Z|\mathbb Z}_{\rho,\sigma}\times_{\gamma_{\rho,\sigma}}N_{\rho}^{*}\times_{\gamma_{\rho,\sigma}}N_{\sigma}^{*}\arrow[r]&S\mathcal{G}_{\rho,\sigma}^{\mathrm{eff}}\\
x=\unsure{\{u,a,b\}|(u,a,b)}\arrow[r,mapsto] & (\alpha,\beta,\mu)
\end{tikzcd}}\right.
\quad\text{avec}\quad
\left\{\begin{array}{l}
\alpha=u\overbrace{a^{-1}\tau'^{\unsure{\nu'(x)}}}^{\in Z_{\rho}}\\
\beta=u\underbrace{b^{-1}\tau'^{\nu(x)}}_{\in Z_{\sigma}}\\
\mu=\chi_{S\Gamma}(x)
\end{array}\right.
\]

\margin{pour être plus correct, il faudrait écrire
\[\left\{\begin{array}{l}
\alpha=u\,i_{\rho}(a)^{-1}\tau'^{\nu'(x)}\\
\beta=u\,i_{\sigma}(b)^{-1}\tau'^{\nu(x)}
\end{array}\right.\]}

où on a défini $\nu'(x),\unsure{\nu(x)|\nu'(x)}\in\mathbb Z/2\mathbb Z$ par
\[
\deleted{\unsure{(}}\qquad
(-1)^{\nu'(x)}=\varepsilon'(x),\qquad(-1)^{\nu(x)}=\varepsilon(x).
\]

Par définition de $S_{0}\Gamma_{\rho,\sigma}$ comme produit fibré, on aura en effet $ua^{-1},ub^{-1}\in\mathcal{G}^{+}$, donc $\alpha,\beta\in\mathcal{G}$, on aura d'ailleurs
\[
\text{(77)}\qquad
\left\{\begin{array}{l}
\displaystyle\alpha(\rho)=u\left(\underbrace{a^{-1}\left(\underbrace{\tau'^{\nu'(x)}(\rho)}_{\rho^{\varepsilon'(x)}}\right)}_{(\rho^{\varepsilon'(x)})^{\varepsilon'(x)}=\rho}\right)=u(\rho),\quad\text{et de même}\\[2em]
\beta(\sigma)=u(\sigma),
\end{array}\right.
\qquad\text{i.e.}\quad
\begin{array}{l}
[\alpha,\rho]=[u,\rho]\\
[\beta,\sigma]=[u,\sigma]
\end{array}
\]

\margin{\uline{NB} On a, pour \unsure{tt} $a\in N_{\rho}^{*}$, si $\varepsilon_{\rho}(a)=(-1)^{\nu_{\rho}(a)}$, $a\tau'^{\nu_{\rho}(a)}\in Z_{\rho}$, donc $[ua\tau'^{\nu_{\rho}(a)},\rho]=[u,\rho]$. Ainsi pour $\sigma\ldots$}

\deleted{de sorte que} d'ailleurs la condition
$\underbrace{[\alpha,\rho]}_{[u,\rho]}=\underbrace{[\beta,\sigma]}_{[u,\sigma]}\varepsilon_{1}^{\mu-1}$ \deleted{\unsure{ici}}
ne fait qu'exprimer la condition que $u(\varepsilon_{0})=\varepsilon_{0}^{\mu}$ qui intervient dans la définition $\overline{\mathrm m}$ de $Z_{\rho,\sigma}$.

\deleted{D'ailleurs} \inserted{Je dis que} l'application envisagée \inserted{(76)} est injective, \deleted{car} Il suffit de voir que quand $(\alpha,\beta,\mu)$ sont connus, on connaît $u$ --- \inserted{en effet} on \deleted{a \unsure{alors}} \inserted{en termes de $(\alpha,\beta,\mu)$}
\[
\text{(78)}\qquad
\left\{\begin{array}{l}
\nu'(x)=\unsure{\nu|\varepsilon}_{\mathcal{G}}(\alpha),\quad
\nu(x)=\unsure{\nu|\varepsilon}_{\mathcal{G}}(\beta),\quad
\deleted{\unsure{\mu=}}\ \chi_{S\Gamma}(x)=\unsure{\chi}_{Z}(u)=\mu\\[0.6em]
\deleted{\text{d'où on tire}}\quad\text{d'où}\quad
a=\tau'^{\unsure{\nu|\varepsilon}_{\mathcal{G}}(\alpha)}\alpha^{-1}u,\quad
b=\tau'^{\unsure{\nu|\varepsilon}_{\mathcal{G}}(\beta)}\beta^{-1}u
\end{array}\right.
\]
\deleted{\unsure{on}} donc si $u$ est connu, on connaît également $a$ et $b$ donc $x$. D'autre part pour déterminer \inserted{en termes de $\alpha,\beta,\mu$,} $u\in Z_{\rho,\sigma}\subset\operatorname{Aut}(\mathcal{G}^{+})\times\widehat{\mathbb Z}^{*}\times\{\pm1\}\times\{\pm1\}$ il suffit, en outre de (78), de \deleted{\unsure{le}} déterminer \deleted{\unsure{en tant}} l'élément de $\operatorname{Aut}(\mathcal{G}^{+})$ qu'il définit, ce qui est

\mspage{566}

{\small
\begin{center}
566\\
\texttt{https://grothendieck.umontpellier.fr}\\
Université de Montpellier, UMR CNRS 5149, Case courrier 051, Place Eugène Bataillon\\
34095 Montpellier cedex 5 - France
\end{center}
}
